\section{First-order unification}
\label{sec:automated-reasoning-03}

The definitions of finite-support simultaneous term substitution, composition,
the instance preorder, unifier, and most general unifier are exactly those in
the first-order resolution chapter.  In this chapter only, substitutions are
written \(\theta_{\mathrm u},\rho_{\mathrm u},\eta_{\mathrm u}\).  The
subscript is necessary because \(\theta\) already denotes a formula in the
preceding Tseitin example; it changes no definition.  Composition is oriented
by
\[
  (\theta_{\mathrm u}\eta_{\mathrm u})(x)
  =(\theta_{\mathrm u}(x))\eta_{\mathrm u}.
\]
Thus \(\theta_{\mathrm u}\) is a most general unifier of atoms \(A\) and
\(B\) iff \(A\theta_{\mathrm u}=B\theta_{\mathrm u}\) literally and every
unifier \(\rho_{\mathrm u}\) has the form
\(\theta_{\mathrm u}\eta_{\mathrm u}\) for some term substitution
\(\eta_{\mathrm u}\).

The unification procedure performs the occurs check when it considers an
equation \(x\doteq t\).  If \(x\) occurs properly in \(t\), the procedure
fails.  Thus the atoms \(P(x)\) and \(P(f(x))\) have no unifier: their
argument equation is \(x\doteq f(x)\).  Assigning \(x\mapsto f(x)\) does not
unify the atoms because their simultaneously substituted arguments are
\(f(x)\) and \(f(f(x))\).

\begin{example}
The atoms \(P(f(x,g(a)))\) and \(P(f(g(y),g(y)))\) have the most general
unifier
\(\theta_{\mathrm u}=\set{x\mapsto g(a),y\mapsto a}\).  Applying it to
either atom yields \(P(f(g(a),g(a)))\).
\end{example}

